\section{Why can’t a Woman be like a Man?}

``Why Can't a Woman Be More Like a Man?'' Rex Harrison famously asked in My Fair Lady\footnote{\url{https://www.youtube.com/watch?v=Doz5w2W-jAY}}. Of course, implicit in that plaint is the assumption that men are ``real'' and women are at best inferior copies. One of my favorite writers on psychology is \textbf{James Hillman} who addresses this topic in his book The Myth of Analysis. Jung says, and it seems to me that Gornahoor is in accord with this, that the spirit is associated with the masculine and matter with the feminine. That is why the worth of women is so often tied to her body, while a man keeps his respect even in his bodily decline; \textbf{Carl Jung} calls her human body, ``the thing most prone to gross material corruption.''

Hillman brings in a number of fascinating ideas. For example, \textbf{Jacob Boehme} said the Eve is Adam's ``sleep''. That is, she is his unconscious. \textbf{Berdyaev} adds that woman brings sexuality into the world. I think this means that sexuality is man's unconscious. It drives him in ways he cannot understand; it embarrasses him when it doesn't frighten him. Women can display their sexuality overtly, they can dress provocatively, they can act sensually. That is because they are sex. Men, on the other hand, are not. They dress in suits that hide their sexuality, which, in any case, is something added to them, but is not their essence. Only homosexuals can flaunt their sexuality like women; straight men who do so---and it is becoming more common with the metrosexuals---are considered to be weak.

So this leads us to the method to make a woman more like a man: it is to deny them as sexual and material beings. A generation ago, \textbf{William Masters} and \textbf{Virginia Johnson} published their ground-breaking book, Human Sexual Response. While this study's aim was to liberate women as sexual beings, Hillman has an interesting take on it.

\begin{quotex}
The emphasis upon female orgasm and upon the ``freedom'' of women through contraception and abortion are the ways in which the dogma of female inferiority is presenting itself at the moment: the prototype for free and healthy sexuality is male; by a technology of the orgasm, by legalizing abortion, and by perfecting the oral contraceptive, women can more closely approximate male sexual patterns.
\end{quotex}


In all the time since its publication in the 60s, no one doubts this emphasis; those who do so are attacked and shamed. However, no one sees this a ``female inferiority'', but rather as the road to equality.

\paragraph{Lies and White Lies}


Nevertheless, even if material and sexual equality has been achieved, this does not mean psychological parity. The clearest expressions of this was in \textbf{Otto Weininger}'s Sex and Character\footnote{\url{http://www.gornahoor.net/library/SexAndCharacter.pdf}}. He claims that the more a woman is like the absolute woman, i.e., the less she is like a man, the more logic and ethics are completely alien to her. It then follows that the woman has no soul nor any possibility of genius. In an obvious allusion to Adam and Eve, Weininger says, ``The man of genius possesses the complete female within himself, but woman herself is only a part of the Universe.''

Weininger dismisses woman's consciousness as consisting of ``henids'', i.e., vague, unconscious feelings. He privileges the allegedly male's consciousness of clear and distinct ideas. But, really, can the world truly be so clear and distinct? A woman feels her way into a situation, she senses things, she notices subtle changes that elude the consciousness of men. Is that not true, is that not valuable? Why are men so fearful of the dark, moist part of the mind?

Weininger claims that women are incapable of knowing the truth, hence they lie. I don't believe that, since it assumes a rather constricted notion of ``truth''. If a boy asks me out for a date, and I tell him I have to take care of my elderly aunt in Hialeah, is that really a lie? He knows what I mean. How would it be better if I said, ``I can't go out with you because I would be embarrassed to be seen in public with you in case one of my friends saw us together.'' The feelings have their own way of knowing the truth.


\flrightit{Posted on 2012-10-27 by Teresa Chicalinda}

\begin{center}* * *\end{center}

\begin{footnotesize}
\begin{sffamily}
\texttt{Bo on 2012-10-28 at 14:24 said:}

The truth is preferable to having to go to Hialeah, good point.

\hfill

\texttt{Cologero on 2012-10-29 at 16:33 said:}

Señorita Chicalinda, you can't judge Tradition by its perverters, who are often misogynists for the reasons you point out. In an initiatic text, The Romance of the Rose, we read from the sermon of Genius (if you catch the irony in comparison to Weininger):

\begin{quote}\footnotesize

When they are first created, God has the same love for all, and gives rational souls to men as well as to women; therefore I believe that he wants every soul, not just one, to follow the best path and to come as quickly as possible to himself.
\end{quote}


\hfill

\texttt{Kenneth Vincent DeWeese on 2012-10-29 at 21:49 said:}

Edith Stein wrote about emotions as providing an intuitive grasp of truth and I think that's accurate. Of course, they can be mistaken and often are, but the intellect is also easily led astray. That's all part of the noetic darkness we inherit in our fallen condition.

\hfill

\texttt{Graham on 2012-10-31 at 00:18 said:}

``If a boy asks me out for a date, and I tell him I have to take care of my elderly aunt in Hialeah, is that really a lie? He knows what I mean. ... The feelings have their own way of knowing the truth.''

I can't argue with the example given, but I will point out that it only shows that women are capable of knowing their own feelings. Which is facile, and isn't what Weininger was talking about here, whatever value we may place on his thought in general. A particular woman's feelings towards a particular man have little to do with the truth in a metaphysical sense, which is what we're interested in -- or even in a philosophical sense, unless we are talking about the philosophy of subjectivism.

And in fact, the man knows what she means because his reason apprehends that she is telling him a lie, not because his feelings do.

To leap from that example, to saying that her feelings know *the* truth, is completely misguided. And I think it demonstrates how many people, and especially modern women, have an extremely high opinion of themselves. To the extent that they will even confuse their feelings with reality itself, such that their experience of their own feelings becomes ``knowledge of the truth.''

I do agree with what Kenneth Vincent DeWeese wrote, however. Emotion can have a reductive, or better put a synthetic power that helps in making judgments quickly, assuming the emotive faculty has been well-formed according to solid principles. And I agree that thought should not be confused with reality either.


\hfill

\end{sffamily}
\end{footnotesize}

